\documentclass[border=4pt]{standalone}
\usepackage{amsmath,amssymb,mathtools,xcolor,varwidth}
\begin{document}
\begin{varwidth}{28em}
\large\bfseries
Draw the tangent at $P$ and the conjugate diameters $DK$ and $PG$: a      pair through the centre $C$ where each runs parallel to the tangent at the end      of the other. Let $Q$ be a body-point near $P$, with the ordinate $Qv$ dropped      to the diameter $PG$, and complete the parallelogram $QxPR$. This is      the ellipse's version of the $QR$ construction of Prop. VI, whose nascent limit      measures the deflection. By Lemma VII the chord $SP$ and its tangent      coincide in the ultimate ratio, licensing this passage to the limit.
\end{varwidth}
\end{document}
